% from ega3-1-fr.tex
\oldpage[III]{148}

$Y_n$-préschéma fini $X_n=\operatorname{Spec}(\mathcal A_n)$ (II, 6.1.3)~;
pour $m\leq n$, l'homomorphisme canonique surjectif
$h_{mn}:\mathcal A_n\to\mathcal A_m$ définit un morphisme
$u_{mn}:X_m\to X_n$ tel que le diagramme
\[
\xymatrix{
X_n\ar[d]_{f_n} & X_m\ar[l]_{u_{mn}}\ar[d]^{f_m}\\
Y_n & Y_m\ar[l]
}
\]
($f_n$ étant le morphisme structural) soit commutatif et identifie $X_m$ au
produit $X_n\times_{Y_n}Y_m$, comme on le voit aussitôt (II, 1.4.6). Le
préschéma formel $\mathfrak X$, limite inductive du système inductif $(X_n)$
est alors localement noethérien et tel que le morphisme structural
$f:\mathfrak X\to\mathfrak Y$, limite inductive du système $(f_n)$, soit fini
(\hyperref[III.4.8.1-fr]{4.8.1} et II, 10.12.3.1)~; on a vu en outre dans la
démonstration de (\hyperref[III.4.8.6-fr]{4.8.6}) que
$f_*(\mathcal O_{\mathfrak X})$ est limite projective des $\mathcal A_n$, donc
égale à $\mathcal A$ (I, 10.10.6). Quant à l'assertion d'unicité, elle est
conséquence du résultat plus général suivant~:
\end{proof}

\begin{proposition}[4.8.8]
\label{III.4.8.8-fr}
Soient $\mathfrak Y$ un préschéma formel localement noethérien,
$\mathfrak X$, $\mathfrak X'$ deux $\mathfrak Y$-préschémas formels finis
au-dessus de $\mathfrak Y$, $f:\mathfrak X\to\mathfrak Y$,
$f':\mathfrak X'\to\mathfrak Y$ les morphismes structuraux. Il existe une
bijection canonique de
\[
\operatorname{Hom}_{\mathfrak Y}(\mathfrak X,\mathfrak X')
\quad\hbox{sur}\quad
\operatorname{Hom}_{\mathcal O_{\mathfrak Y}}
\bigl(f'_*(\mathcal O_{\mathfrak X'}),f_*(\mathcal O_{\mathfrak X})\bigr)
\footnotemark.
\]
\footnotetext{La dernière expression désigne l'ensemble des homomorphismes de
$\mathcal O_{\mathfrak Y}$-Algèbres
$f'_*(\mathcal O_{\mathfrak X'})\to f_*(\mathcal O_{\mathfrak X})$.}
\end{proposition}

\begin{proof}
La définition de cette application $h\to\mathcal A(h)$ est la même que dans
(II, 1.1.2), et pour voir qu'elle est bijective, on est aussitôt ramené au cas
où $\mathfrak Y=\operatorname{Spf}(B)$ est un schéma formel affine noethérien.
Mais alors $\mathfrak X=\operatorname{Spf}(A)$,
$\mathfrak X'=\operatorname{Spf}(A')$, où $A$ et $A'$ sont deux $B$-algèbres
finies et $f_*(\mathcal O_{\mathfrak X})=A^\Delta$,
$f'_*(\mathcal O_{\mathfrak X'})=A'^\Delta$. La conclusion résulte alors de
la correspondance biunivoque, d'une part entre les $\mathfrak Y$-morphismes
$\mathfrak X\to\mathfrak X'$ et les $B$-homomorphismes (nécessairement
continus) $A'\to A$ qui sont des homomorphismes d'algèbres (I, 10.2.2), et
d'autre part entre les homomorphismes de $B$-modules $A'\to A$ et les
homomorphismes de $\mathcal O_{\mathfrak Y}$-Modules
$A'^\Delta\to A^\Delta$ (I, 10.10.2.3).
\end{proof}

\begin{corollary}[4.8.9]
\label{III.4.8.9-fr}
Dans la correspondance biunivoque canonique définie dans
(\hyperref[III.4.8.8-fr]{4.8.8}), les immersions fermées
$\mathfrak X\to\mathfrak X'$ correspondent aux homomorphismes surjectifs de
$\mathcal O_{\mathfrak Y}$-Algèbres
$f'_*(\mathcal O_{\mathfrak X'})\to f_*(\mathcal O_{\mathfrak X})$.
\end{corollary}

\begin{proof}
La question étant encore locale sur $\mathfrak Y$, on est ramené à la
définition des immersions fermées de préschémas formels localement noethériens
(I, 10.14.2).
\end{proof}

\begin{corollary}[4.8.10]
\label{III.4.8.10-fr}
Les notations et hypothèses étant celles de
(\hyperref[III.4.8.1-fr]{4.8.1}), pour qu'un morphisme adique $f$ soit une
immersion fermée, il faut et il suffit que $f_0$ soit une immersion fermée (de
préschémas usuels).
\end{corollary}

\begin{proof}
Cela résulte aussitôt de (\hyperref[III.4.8.9-fr]{4.8.9}) et de la condition de
surjectivité pour un homomorphisme de $\mathcal O_{\mathfrak Y}$-Modules
cohérents (I, 10.11.5).
\end{proof}

\begin{proposition}[4.8.11]
\label{III.4.8.11-fr}
Pour qu'un morphisme $f:\mathfrak X\to\mathfrak Y$ de préschémas formels localement

